\documentclass[11pt,twoside,openright]{book}

% =============================================================================
% Mass Spectrometry Data Analysis with R
% Chapter 09: Metabolite Annotation and Adduct Deconvolution
% Author: Lucas VHH TRAN (CK-CARE Davos)
% Style: Springer Nature Monograph (Health Analytics with R Pedagogy)
% =============================================================================

\usepackage{springer_monograph}
\usepackage{placeins}
\usepackage{float}

\hypersetup{
    pdftitle={Chapter 09: Metabolite Annotation and Adduct Deconvolution},
    pdfauthor={Lucas VHH TRAN}
}

\begin{document}

% --- Chapter Title Block ---
\setcounter{chapter}{8}
\chapter{Metabolite Annotation and Adduct Deconvolution}
\label{ch:metabolite_annotation}

\noindent
\textbf{Lucas VHH TRAN}\footnote{CK-CARE Davos, Christine K\"uhne--Center for Allergy Research and Education, Davos, Switzerland. Email: \texttt{lucas.tran@ck-care.ch}}\\
\rule{\textwidth}{0.8pt}

\begin{abstract}
\noindent
Feature detection extracts thousands of chromatographic peaks defined by mass-to-charge ($m/z$) and retention time ($t_R$), but a coordinate pair does not identify a molecule. In electrospray ionization (ESI), single chemical metabolites partition into multiple spectral features: protonated ions, sodium and potassium adducts, neutral water losses, and isotopic envelopes. Conversely, a single accurate mass matches multiple isomeric structures. Metabolite annotation deconvolutes redundant adduct multiplets into discrete parent molecules \cite{kuhl2012camera, rainer2022modular} and queries reference chemical databases \cite{wishart2022hmdb, guijas2018metlin, horai2010massbank} within instrument-calibrated mass tolerances. This chapter establishes the physical chemistry, mathematical formulations, and Bioconductor implementations for metabolite annotation. We derive exact monoisotopic mass arithmetic from atomic mass defects \cite{gross2017mass}, formulate tolerance-aware search intervals in parts-per-million (ppm) \cite{kind2007seven}, and model natural isotopic fine structure distributions using multinomial expansions \cite{duehrkop2019sirius}. We present the Isotope Envelope Discriminability Theorem, demonstrating the resolving power required to separate sulfur-34 ($^{34}\text{S}$) from double carbon-13 ($^{13}\text{C}_2$) substitutions ($10.91$ mDa). Using open access benchmark data, we execute a modernized Bioconductor workflow in R 4.6.1 combining \texttt{MetaboCoreUtils}, \texttt{MetaboAnnotation}, and \texttt{SummarizedExperiment} \cite{rainer2022modular, huber2015orchestrating, smith2006xcms, tautenhahn2008centwave}. We demonstrate adduct deconvolution, tolerance-controlled database matching, and structured reporting conforming to the Metabolomics Standards Initiative (MSI) confidence tiers \cite{sumner2007proposed, schymanski2014identifying}. The chapter concludes with the Five Deadly Sins of metabolite annotation, a diagnostic checklist, and four multi-tier exercises with complete mathematical derivations and verified R code.
\end{abstract}

\noindent
\textbf{Keywords:}\enspace Metabolite Annotation, Adduct Deconvolution, Mass Accuracy, Isotopic Fine Structure, Formula Assignment, \texttt{MetaboCoreUtils}, ESI-MS.

\vspace{0.8em}
\noindent\rule{\textwidth}{0.4pt}

% =============================================================================
% Methodological Roadmap
% =============================================================================

\begin{table}[ht]
\centering
\caption{Methodological Roadmap: Four Stages of Metabolite Annotation}
\label{tab:annotation_roadmap}
\small
\begin{tabularx}{\textwidth}{l p{3.8cm} p{4.2cm} X}
\toprule
\textbf{Stage} & \textbf{Statistical Mechanism} & \textbf{Key Analytical Objective} & \textbf{Primary Failure Mode \& Remedy} \\
\midrule
\textbf{1. Deconvolution} & Peak shape correlation ($r \ge 0.85$) and chromatographic grouping \cite{kuhl2012camera, pluskal2010mzmine}. & Collapse co-eluting adducts and isotopes into pseudo-spectra. & Treating adducts as distinct biological features. Remedy: Group by retention time and EIC shape correlation. \\
\addlinespace
\textbf{2. Adduct Inversion} & Algebraic neutral mass calculation from known adducts \cite{gross2017mass, rainer2022modular}. & Deduce monoisotopic neutral precursor mass $M_0$. & Inverting with wrong adduct assumption (e.g., Na vs H). Remedy: Multi-adduct consistency verification. \\
\addlinespace
\textbf{3. Exact Search} & Symmetric parts-per-million windowing ($\pm \delta_{\text{ppm}}$) \cite{kind2007seven, wishart2022hmdb}. & Query reference databases for candidate formulas and structures. & Search space explosion from wide tolerance ($> 15$ ppm). Remedy: Instrument lock-mass calibration ($\le 5$ ppm). \\
\addlinespace
\textbf{4. Fine Structure} & Multinomial isotopic envelope simulation and cosine scoring \cite{duehrkop2019sirius, tsugawa2015msdial}. & Filter candidate formulas via $M+1$ and $M+2$ abundance matching. & Formula degeneracy above $m/z = 300$. Remedy: Score fine-structure mass shifts ($^{34}\text{S}$ vs $^{13}\text{C}_2$). \\
\bottomrule
\end{tabularx}
\end{table}

\FloatBarrier

\begin{notebox}
\textbf{Tutor's Methodological Advice: The Anomaly of the Accurate Mass}\\
An accurate mass match provides an elemental formula, not a confirmed chemical entity. Above $m/z = 300$, dozens of distinct chemical structures share the identical molecular formula. Reporting a database match as a confirmed metabolite without matching retention times against an authentic chemical standard or matching MS/MS fragmentation against a reference spectral library violates reporting standards. Always assign and report annotations according to the formal Metabolomics Standards Initiative (MSI) confidence levels.
\end{notebox}

% =============================================================================
\section{Biophysical and Mathematical Foundations}
% =============================================================================

\subsection{Electrospray Ionization Thermodynamics and Adduct Chemistry}

Electrospray ionization (ESI) transfers neutral molecules ($M$) dissolved in a liquid mobile phase into gas-phase ions inside the mass spectrometer. Analytes acquire electrical charge through proton transfer or coordination with ubiquitous cations and anions in solution \cite{gross2017mass}.

In positive ionization mode ($\text{ESI}^+$), basic nitrogen atoms and polar oxygen functional groups acquire protons or coordinate alkali metals:
\begin{align}
\text{Protonated molecule: } & M + \text{H}^+ \xrightarrow{} [M+\text{H}]^+, & (m/z) &= M_0 + 1.007276 \\
\text{Sodium adduct: } & M + \text{Na}^+ \xrightarrow{} [M+\text{Na}]^+, & (m/z) &= M_0 + 22.989218 \\
\text{Potassium adduct: } & M + \text{K}^+ \xrightarrow{} [M+\text{K}]^+, & (m/z) &= M_0 + 38.963158 \\
\text{Ammonium adduct: } & M + \text{NH}_4^+ \xrightarrow{} [M+\text{NH}_4]^+, & (m/z) &= M_0 + 18.033826 \\
\text{In-source water loss: } & [M+\text{H}]^+ - \text{H}_2\text{O} \xrightarrow{} [M+\text{H}-\text{H}_2\text{O}]^+, & (m/z) &= M_0 - 17.003289
\end{align}
where $M_0$ represents the monoisotopic neutral mass of the analyte, and mass values are corrected for electron mass ($m_e = 0.00054858\;\text{Da}$).

In negative ionization mode ($\text{ESI}^-$), acidic functional groups (carboxylic acids, phosphates, sulfates) deprotonate or bind mobile-phase acid counterions:
\begin{align}
\text{Deprotonated molecule: } & M - \text{H}^+ \xrightarrow{} [M-\text{H}]^-, & (m/z) &= M_0 - 1.007276 \\
\text{Formate adduct: } & M + \text{HCOO}^- \xrightarrow{} [M+\text{HCOO}]^-, & (m/z) &= M_0 + 44.998201 \\
\text{Acetate adduct: } & M + \text{CH}_3\text{COO}^- \xrightarrow{} [M+\text{CH}_3\text{COO}]^-, & (m/z) &= M_0 + 59.013851
\end{align}

If an analyst treats each detected peak as a distinct biological molecule, a single metabolite generates five to eight redundant features. If that metabolite changes across experimental conditions, all eight features are declared statistically significant in differential testing, distorting false discovery rates and pathway analysis \cite{kuhl2012camera}.

\subsection{Exact Monoisotopic Mass Arithmetic and Atomic Mass Defects}

The unified atomic mass unit (Dalton, $\text{Da}$) is defined such that one atom of carbon-12 ($^{12}\text{C}$) in its nuclear ground state has a mass of exactly $12.0000000\;\text{Da}$. All other atomic nuclei exhibit non-zero **mass defects** resulting from nuclear binding energy:
\begin{equation}
\Delta m_{\text{defect}} = m_{\text{exact}} - A
\end{equation}
where $A$ represents the nominal mass number (total nucleon count):
\begin{itemize}
    \item Hydrogen ($^{1}\text{H}$): $m = 1.0078250\;\text{Da}$ ($\Delta m = +0.007825\;\text{Da}$, positive defect).
    \item Nitrogen ($^{14}\text{N}$): $m = 14.0030740\;\text{Da}$ ($\Delta m = +0.003074\;\text{Da}$, positive defect).
    \item Oxygen ($^{16}\text{O}$): $m = 15.9949146\;\text{Da}$ ($\Delta m = -0.005085\;\text{Da}$, negative defect).
    \item Phosphorus ($^{31}\text{P}$): $m = 30.9737620\;\text{Da}$ ($\Delta m = -0.026238\;\text{Da}$, negative defect).
    \item Sulfur ($^{32}\text{S}$): $m = 31.9720707\;\text{Da}$ ($\Delta m = -0.027929\;\text{Da}$, negative defect).
\end{itemize}

Let $\mathcal{E} = \{\text{C}, \text{H}, \text{N}, \text{O}, \text{P}, \text{S}\}$ denote the primary bio-elements. The exact neutral monoisotopic mass $M_0$ of a molecule with formula $\text{C}_{n_{\text{C}}}\text{H}_{n_{\text{H}}}\text{N}_{n_{\text{N}}}\text{O}_{n_{\text{O}}}\text{P}_{n_{\text{P}}}\text{S}_{n_{\text{S}}}$ is:
\begin{equation}
M_0 = \sum_{e \in \mathcal{E}} n_e \, m_e
\label{eq:neutral_mass_exact}
\end{equation}

\subsection{Tolerance-Aware Search Geometry and Parts-Per-Million Metrics}

Because mass spectrometers measure with finite physical precision, database search algorithms match observed peaks against theoretical targets within an instrument-specific tolerance window.

The relative mass error $\delta_{\text{ppm}}$ in parts-per-million (ppm) is:
\begin{equation}
\delta_{\text{ppm}} = \frac{m_{\text{obs}} - m_{\text{theo}}}{m_{\text{theo}}} \times 10^6
\label{eq:ppm_error}
\end{equation}
A symmetric tolerance window of $\pm \tau_{\text{ppm}}$ around an observed peak $m_{\text{obs}}$ defines the search interval:
\begin{equation}
\left[ m_{\text{obs}} \left(1 - \frac{\tau_{\text{ppm}}}{10^6}\right), \; m_{\text{obs}} \left(1 + \frac{\tau_{\text{ppm}}}{10^6}\right) \right]
\label{eq:search_interval}
\end{equation}
On modern Orbitrap and Q-TOF mass spectrometers, search tolerances are constrained to $\tau_{\text{ppm}} \le 5.0\;\text{ppm}$. Applying wider windows ($>15$ ppm) causes search space explosion, returning dozens of false positive candidates for each query feature \cite{kind2007seven}.

\subsection{Isotopic Fine Structure and Envelope Discriminability}

Natural elements exist as mixtures of stable isotopes with known natural abundances:
\begin{itemize}
    \item Carbon: $^{12}\text{C}$ ($98.93\%$), $^{13}\text{C}$ ($1.07\%$, $\Delta m = +1.003355\;\text{Da}$).
    \item Nitrogen: $^{14}\text{N}$ ($99.63\%$), $^{15}\text{N}$ ($0.37\%$, $\Delta m = +0.997035\;\text{Da}$).
    \item Oxygen: $^{16}\text{O}$ ($99.76\%$), $^{17}\text{O}$ ($0.038\%$, $\Delta m = +1.004217\;\text{Da}$), $^{18}\text{O}$ ($0.205\%$, $\Delta m = +2.004245\;\text{Da}$).
    \item Sulfur: $^{32}\text{S}$ ($94.99\%$), $^{33}\text{S}$ ($0.75\%$, $\Delta m = +0.999388\;\text{Da}$), $^{34}\text{S}$ ($4.21\%$, $\Delta m = +1.995796\;\text{Da}$).
\end{itemize}

The relative abundance of the $M+1$ isotopic peak relative to the monoisotopic peak $M$ is governed by binomial probabilities:
\begin{equation}
\frac{I(M+1)}{I(M)} \approx n_{\text{C}} \times 0.0107 + n_{\text{N}} \times 0.0037 + n_{\text{H}} \times 0.000115
\label{eq:m1_ratio}
\end{equation}

At $M+2$, different elemental substitutions produce distinct exact masses. High-resolution mass spectrometry resolves these isotopic fine structures:

\vspace{0.6em}
\noindent\textbf{Theorem 9.1 (Isotope Envelope Discriminability).}\enspace\textit{Let a molecule incorporate either a single sulfur-34 ($^{34}\text{S}$) substitution or a double carbon-13 ($^{13}\text{C}_2$) substitution, both yielding nominal mass shift $+2\;\text{Da}$. The physical mass separation between the two isotopic peaks is:
\begin{equation}
\Delta m(^{34}\text{S} - {}^{13}\text{C}_2) = |1.995796 - 2.006709| = 0.010913\;\text{Da} \approx 10.91\;\text{mDa}
\end{equation}
To achieve baseline mass resolution ($R_{\text{FWHM}} = m / \Delta m$) between these two isotopic peaks at precursor mass $m/z = 257$, the mass spectrometer must operate at a resolving power satisfying:
\begin{equation}
R \ge \frac{257}{0.010913} \approx 23,550
\end{equation}}

\vspace{0.4em}
\noindent\textbf{Proof.}\enspace The exact mass shift for a $^{34}\text{S}$ substitution is:
\begin{equation*}
\Delta m(^{34}\text{S}) = 33.967867 - 31.972071 = 1.995796\;\text{Da}
\end{equation*}
The exact mass shift for two $^{13}\text{C}$ substitutions is:
\begin{equation*}
\Delta m(^{13}\text{C}_2) = 2 \times 1.0033548 = 2.006710\;\text{Da}
\end{equation*}
The physical mass separation is $\Delta m = |2.006710 - 1.995796| = 0.010914\;\text{Da} = 10.914\;\text{mDa}$. Baseline resolution requires peak full-width at half-maximum $\text{FWHM} \le \Delta m$. Thus, $R = m / \text{FWHM} \ge 257 / 0.010914 \approx 23,548$. Instruments operating at $R \ge 30,000$ (e.g., Orbitrap or high-resolution Q-TOF) resolve this doublet cleanly. \hfill $\blacksquare$

\vspace{0.6em}
The similarity between an experimental isotopic envelope $\mathbf{x} = (x_0, x_1, x_2)$ and a theoretical candidate $\mathbf{y} = (y_0, y_1, y_2)$ is scored using the normalized dot product:
\begin{equation}
S_{\text{iso}} = \frac{\sum_{i=0}^k x_i \, y_i}{\sqrt{\sum_{i=0}^k x_i^2} \, \sqrt{\sum_{i=0}^k y_i^2}}
\label{eq:iso_score}
\end{equation}
Candidates with $S_{\text{iso}} < 0.90$ are rejected as chemically implausible \cite{duehrkop2019sirius, tsugawa2015msdial}.

% =============================================================================
\section{Algorithmic Architecture and Container Engineering}
% =============================================================================

\subsection{Peak Shape Correlation Deconvolution}

Deconvolution partitions raw feature tables into pseudo-spectra by exploiting chromatographic co-elution and peak shape correlation across scans. If two features $A$ and $B$ originate from the same chemical metabolite (e.g., $[M+\text{H}]^+$ and its $[M+\text{Na}]^+$ adduct), their extracted ion chromatogram (EIC) intensities across retention time scans $t_1, t_2, \dots, t_S$ exhibit strong linear correlation:
\begin{equation}
r_{AB} = \frac{\sum_{s=1}^S (I_{A, s} - \bar{I}_A)(I_{B, s} - \bar{I}_B)}{\sqrt{\sum_{s=1}^S (I_{A, s} - \bar{I}_A)^2} \, \sqrt{\sum_{s=1}^S (I_{B, s} - \bar{I}_B)^2}} \ge 0.85
\label{eq:eic_cor}
\end{equation}
Features satisfying $r_{AB} \ge 0.85$ and apex retention time alignment $\Delta t_R \le 0.02$ min are grouped into a single pseudo-spectrum. Pairwise mass differences within the group are matched against an adduct library (Table~\ref{tab:annotation_roadmap}) to identify the monoisotopic precursor $[M+\text{H}]^+$ \cite{kuhl2012camera, pluskal2010mzmine}.

\subsection{Bioconductor Modernization: MetaboCoreUtils and MetaboAnnotation}

Legacy metabolomics pipelines relied on \texttt{CAMERA}, which operated on old \texttt{xcmsSet} data structures and required complex S4 conversions. Modern Bioconductor workflows decouple formula arithmetic, adduct transformation, and spectral matching into dedicated modular packages \cite{rainer2022modular, huber2015orchestrating, smith2006xcms, tautenhahn2008centwave}:
\begin{itemize}
    \item \texttt{MetaboCoreUtils}: C-level vectorized formula calculations, atomic mass arithmetic, adduct parsing (\texttt{calculateMass}, \texttt{mass2mz}, \texttt{mz2mass}).
    \item \texttt{MetaboAnnotation}: Query-to-reference matching pipelines supporting multi-dimensional constraints (exact mass, retention time, MS/MS spectra).
    \item \texttt{SummarizedExperiment}: The standard rectangular container where feature annotations populate the \texttt{rowData} slot, preserving links to sample metadata (\texttt{colData}) and quantitative abundance matrices (\texttt{assay}).
\end{itemize}

% =============================================================================
\section{Container and Workflow Implementations in R}
% =============================================================================

The end-to-end annotation pipeline proceeds in four computational stages:
\begin{enumerate}
    \item \textbf{Feature Matrix Deconvolution}: Group features by retention time and peak shape correlation ($r \ge 0.85$) to collapse adducts.
    \item \textbf{Neutral Mass Inversion}: Invert observed $[M+\text{Na}]^+$ or $[M+\text{K}]^+$ ions to neutral mass $M_0$ using \texttt{MetaboCoreUtils::mz2mass}.
    \item \textbf{Tolerance-Controlled Search}: Match neutral masses against reference databases (HMDB, KEGG, ChEBI \cite{wishart2022hmdb, guijas2018metlin, horai2010massbank}) within $\pm 5.0\;\text{ppm}$.
    \item \textbf{Isotopic Fine Structure Pruning}: Simulate theoretical isotopic envelopes and filter candidate formulas by cosine score $S_{\text{iso}} \ge 0.90$.
    \item \textbf{Rectangular Packaging}: Integrate candidate annotations, matched formulas, and MSI confidence levels into the \texttt{rowData} of a \texttt{SummarizedExperiment}.
\end{enumerate}

% =============================================================================
\section{Critical Considerations}
% =============================================================================

\subsection{MSI Confidence Levels Quoting Sumner et al. (2007)}

The Metabolomics Standards Initiative (MSI) defines four unambiguous confidence tiers \cite{sumner2007proposed, schymanski2014identifying, broadhurst2018guidelines, wang2016sharing}:
\begin{itemize}
    \item \textbf{Level 1 (Identified Metabolites)}: Confirmed by matching at least two orthogonal properties (retention time, exact mass, and MS/MS fragmentation) against an authentic chemical reference standard analyzed on the identical instrument under identical analytical conditions.
    \item \textbf{Level 2 (Putatively Annotated Compounds)}: Structurally annotated by matching MS/MS fragmentation spectra against public or commercial reference libraries without an authentic standard analyzed locally.
    \item \textbf{Level 3 (Putatively Characterized Compound Classes)}: Assigned to a chemical class based on characteristic physicochemical properties, accurate mass, and diagnostic fragment motifs, while specific constitutional branching or stereochemistry remains unresolved.
    \item \textbf{Level 4 (Unidentified Features)}: Reproducible chromatographic peaks with exact $m/z$ and retention time coordinates that remain unmatched to any database or formula.
\end{itemize}

Full-scan accurate mass matching alone provides **MSI Level 4** (formula level) or **MSI Level 3** (class level). Claiming Level 1 or Level 2 identification from accurate mass alone without MS/MS spectra or authentic standards is a severe scientific reporting defect.

\subsection{The Isomer Dilemma and Orthogonal Dimensions}

Exact mass identifies molecular formulas, not three-dimensional structures. In human serum, hexose sugars illustrate structural degeneracy:
\begin{equation}
\text{Glucose, Fructose, Galactose, and myo-Inositol} \implies \text{Formula: } \text{C}_6\text{H}_{12}\text{O}_6 \quad (M_0 = 180.063388\;\text{Da})
\end{equation}
All four compounds share the identical elemental formula, exact mass, and identical isotopic fine structure. Mass spectrometry alone cannot distinguish them without orthogonal separation dimensions:
\begin{enumerate}
    \item \textbf{Retention Time Matching ($\Delta t_R$)}: Hydrophilic interaction liquid chromatography (HILIC) baseline-resolves glucose, galactose, and fructose.
    \item \textbf{Ion Mobility Spectrometry (IMS)}: Measures gas-phase rotational collision cross sections ($\text{CCS}$ in $\text{\AA}^2$). Structural isomers with different spatial geometries exhibit different drift velocities through buffer gas.
    \item \textbf{Tandem Mass Spectrometry (MS/MS)}: Collision-induced dissociation yields distinct diagnostic fragment ions or cross-ring cleavages under controlled collision energies.
\end{enumerate}

% =============================================================================
\section{Hands-on R Laboratory}
% =============================================================================

The following laboratory scripts execute under R 4.6.1 / Bioconductor 3.20 using \texttt{MetaboCoreUtils} and \texttt{dplyr}. All console outputs represent real executions.

\subsection{Listing 9.1: Exact Mass and Adduct Inversion Arithmetic}

We compute theoretical exact masses for key metabolites and evaluate electrospray adduct transformations in positive mode.

\begin{rchunk}{Listing 9.1: Exact Mass and Adduct Calculation with MetaboCoreUtils}
library(MetaboCoreUtils)
library(dplyr)

# Neutral formulas for key biological metabolites
metabolites <- c(
  Glucose     = "C6H12O6",
  Dopamine    = "C8H11NO2",
  Caffeine    = "C8H10N4O2",
  Tryptophan  = "C11H12N2O2"
)

# Calculate exact neutral monoisotopic masses
exact_masses <- MetaboCoreUtils::calculateMass(metabolites)

# Positive-mode adducts to evaluate
target_adducts_pos <- c("[M+H]+", "[M+Na]+", "[M+K]+", "[M+NH4]+", "[M+H-H2O]+")

# Compute theoretical m/z for Dopamine across adducts
dopamine_mz <- MetaboCoreUtils::mass2mz(exact_masses["Dopamine"], adduct = target_adducts_pos)

adduct_summary <- data.frame(
  Compound       = rep("Dopamine (C8H11NO2)", length(target_adducts_pos)),
  Neutral_Mass   = rep(unname(round(exact_masses["Dopamine"], 4)), length(target_adducts_pos)),
  Adduct         = target_adducts_pos,
  Theoretical_mz = round(as.numeric(dopamine_mz[1, ]), 4),
  stringsAsFactors = FALSE
)

# Demonstrate neutral mass inversion from observed [M+Na]+ ion
obs_mz <- 176.0682   # Observed Dopamine [M+Na]+
inv_mass <- MetaboCoreUtils::mz2mass(obs_mz, adduct = "[M+Na]+")
mass_err_ppm <- ((inv_mass - exact_masses["Dopamine"]) / exact_masses["Dopamine"]) * 1e6

cat("Dopamine Electrospray Adduct Profile (ESI+):\n")
print(adduct_summary)
cat(sprintf("\nAdduct Inversion Audit: Observed [M+Na]+ = %.4f Da\n", obs_mz))
cat(sprintf("  Inferred Neutral Mass: %.4f Da | True Mass: %.4f Da | Delta: %+.4f Da (%+.2f ppm)\n",
            inv_mass, exact_masses["Dopamine"], inv_mass - exact_masses["Dopamine"], mass_err_ppm))
\end{rchunk}

\begin{routput}
Dopamine Electrospray Adduct Profile (ESI+):
             Compound Neutral\_Mass     Adduct Theoretical\_mz
1 Dopamine (C8H11NO2)      153.079     [M+H]+       154.0863
2 Dopamine (C8H11NO2)      153.079    [M+Na]+       176.0682
3 Dopamine (C8H11NO2)      153.079     [M+K]+       192.0421
4 Dopamine (C8H11NO2)      153.079   [M+NH4]+       171.1128
5 Dopamine (C8H11NO2)      153.079 [M+H-H2O]+       136.0757

Adduct Inversion Audit: Observed [M+Na]+ = 176.0682 Da
  Inferred Neutral Mass: 153.0790 Da | True Mass: 153.0790 Da | Delta: +0.0000 Da (+0.01 ppm)
\end{routput}

\subsection{Listing 9.2: Peak Shape Correlation and Multiplet Deconvolution}

We evaluate chromatographic co-elution and peak shape correlation across scans to group adduct multiplets into a single pseudo-spectrum.

\begin{rchunk}{Listing 9.2: Peak Shape Correlation Deconvolution}
# Simulate chromatographic time series across elution window (t = 2.0 to 2.8 min)
time_scans <- seq(2.0, 2.8, by = 0.02)
rt_apex <- 2.40
peak_sigma <- 0.06

# Simulate co-eluting features from a single metabolite:
# FT_181_MH: [M+H]+, FT_163_Loss: [M+H-H2O]+, FT_203_MNa: [M+Na]+
int_mh   <- 120000 * exp(-0.5 * ((time_scans - rt_apex) / peak_sigma)^2) + rnorm(length(time_scans), 0, 800)
int_loss <-  24000 * exp(-0.5 * ((time_scans - rt_apex) / peak_sigma)^2) + rnorm(length(time_scans), 0, 300)
int_mna  <-  48000 * exp(-0.5 * ((time_scans - (rt_apex + 0.005)) / peak_sigma)^2) + rnorm(length(time_scans), 0, 500)
int_cont <-  35000 * exp(-0.5 * ((time_scans - 2.46) / peak_sigma)^2) + rnorm(length(time_scans), 0, 400)

eic_matrix <- cbind(
  FT_181_MH   = pmax(0, int_mh),
  FT_163_Loss = pmax(0, int_loss),
  FT_203_MNa  = pmax(0, int_mna),
  FT_unrel    = pmax(0, int_cont)
)

cor_mat <- round(cor(eic_matrix), 3)
cat("Chromatographic Peak Shape Correlation Matrix (Pearson r):\n")
print(cor_mat)

# Deconvolution grouping threshold: r >= 0.85
cluster_members <- colnames(eic_matrix)[cor_mat["FT_181_MH", ] >= 0.85]
cat(sprintf("\nDeconvoluted Pseudo-Spectrum Group 1 (r >= 0.85 with [M+H]+):\n"))
print(cluster_members)
\end{rchunk}

\begin{routput}
Chromatographic Peak Shape Correlation Matrix (Pearson r):
            FT\_181\_MH FT\_163\_Loss FT\_203\_MNa FT\_unrel
FT\_181\_MH       1.000       0.999      0.998    0.702
FT\_163\_Loss     0.999       1.000      0.998    0.706
FT\_203\_MNa      0.998       0.998      1.000    0.740
FT\_unrel        0.702       0.706      0.740    1.000

Deconvoluted Pseudo-Spectrum Group 1 (r >= 0.85 with [M+H]+):
{}[1] "FT\_181\_MH"   "FT\_163\_Loss" "FT\_203\_MNa" 
\end{routput}

\subsection{Listing 9.3: Tolerance-Aware Database Matching with PPM Error Scoring}

We match neutral query masses against reference chemical databases using symmetric parts-per-million windows.

\begin{rchunk}{Listing 9.3: Tolerance-Aware Database Search}
# Curated reference database snippet
ref_database <- data.frame(
  Compound_Name = c("Glucose", "Fructose", "myo-Inositol", "Dopamine", "Caffeine", "Citric acid"),
  Formula       = c("C6H12O6", "C6H12O6", "C6H12O6", "C8H11NO2", "C8H10N4O2", "C6H8O7"),
  Exact_Mass    = c(180.06339, 180.06339, 180.06339, 153.07898, 194.08038, 192.02700),
  stringsAsFactors = FALSE
)

# Experimental observed features (neutral mass deconvoluted)
query_features <- data.frame(
  Feature_ID     = c("FT_001", "FT_002", "FT_003"),
  Observed_Mass  = c(180.0639, 153.0784, 312.1245),
  Retention_Time = c(2.45, 1.82, 8.40),
  stringsAsFactors = FALSE
)

ppm_tolerance <- 5.0
matches_list <- list()

for (i in seq_len(nrow(query_features))) {
  obs_m <- query_features$Observed_Mass[i]
  delta_limit <- obs_m * (ppm_tolerance / 1e6)
  candidates <- ref_database %>%
    filter(abs(Exact_Mass - obs_m) <= delta_limit)
  
  if (nrow(candidates) > 0) {
    candidates$Feature_ID     <- query_features$Feature_ID[i]
    candidates$Observed_Mass  <- obs_m
    candidates$Retention_Time <- query_features$Retention_Time[i]
    candidates$PPM_Error      <- round(abs(candidates$Exact_Mass - obs_m) / candidates$Exact_Mass * 1e6, 2)
    matches_list[[length(matches_list) + 1]] <- candidates
  }
}

matched_df <- do.call(rbind, matches_list) %>%
  select(Feature_ID, Observed_Mass, Retention_Time, Compound_Name, Formula, PPM_Error)

cat(sprintf("Exact Mass Database Search Results (Tolerance = +/- %.1f ppm):\n", ppm_tolerance))
print(matched_df)
\end{rchunk}

\begin{routput}
Exact Mass Database Search Results (Tolerance = +/- 5.0 ppm):
  Feature\_ID Observed\_Mass Retention\_Time Compound\_Name  Formula PPM\_Error
1     FT\_001      180.0639           2.45       Glucose  C6H12O6      2.83
2     FT\_001      180.0639           2.45      Fructose  C6H12O6      2.83
3     FT\_001      180.0639           2.45  myo-Inositol  C6H12O6      2.83
4     FT\_002      153.0784           1.82      Dopamine C8H11NO2      3.79
\end{routput}

\subsection{Listing 9.4: High-Resolution Isotopic Fine Structure Modeling}

We evaluate the fine mass separation between sulfur-34 and double carbon-13 substitutions at $m/z = 257$, computing required instrument resolving power.

\begin{rchunk}{Listing 9.4: Isotopic Fine Structure Mass Separation}
# Contrast 13C2 double substitution vs 34S sulfur substitution:
# Carbon-13 double substitution: Delta m = 2 * 1.003355 = 2.006710 Da
# Sulfur-34 substitution:        Delta m = 1.995796 Da
fine_mass_diff <- abs(2.006710 - 1.995796)
resolving_power_req <- 257 / fine_mass_diff

cat("High-Resolution Isotopic Fine Structure Mass Differences:\n")
cat(sprintf("  13C2 Isotope Mass Shift:  +%.6f Da\n", 2.006710))
cat(sprintf("  34S  Isotope Mass Shift:  +%.6f Da\n", 1.995796))
cat(sprintf("  Mass Separation (Delta m): %.6f Da (%.2f mDa)\n", fine_mass_diff, fine_mass_diff * 1000))
cat(sprintf("  Resolving Power Required to baseline-resolve 34S from 13C2 at m/z 257: R >= %.0f\n",
            resolving_power_req))
\end{rchunk}

\begin{routput}
High-Resolution Isotopic Fine Structure Mass Differences:
  13C2 Isotope Mass Shift:  +2.006710 Da
  34S  Isotope Mass Shift:  +1.995796 Da
  Mass Separation (Delta m): 0.010914 Da (10.91 mDa)
  Resolving Power Required to baseline-resolve 34S from 13C2 at m/z 257: R >= 23548
\end{routput}

\subsection{Listing 9.5: Structured Annotation Assembly and MSI Level Classification}

We assemble candidate matches into a consolidated annotation report with formal MSI confidence levels.

\begin{rchunk}{Listing 9.5: Annotation Report Assembly}
annotation_report <- query_features %>%
  left_join(matched_df, by = c("Feature_ID", "Observed_Mass", "Retention_Time")) %>%
  group_by(Feature_ID, Observed_Mass, Retention_Time) %>%
  summarise(
    Matched_Count  = sum(!is.na(Compound_Name)),
    Candidates     = if (sum(!is.na(Compound_Name)) > 0) paste(Compound_Name, collapse = " | ") else "Unmatched",
    Formulas       = if (sum(!is.na(Formula)) > 0) paste(unique(Formula), collapse = " | ") else "None",
    Min_PPM_Error  = if (sum(!is.na(PPM_Error)) > 0) min(PPM_Error) else NA_real_,
    .groups        = "drop"
  ) %>%
  mutate(
    MSI_Confidence = case_when(
      Matched_Count == 0 ~ "Level 4 (Unidentified Feature)",
      Matched_Count == 1 ~ "Level 3 (Putatively Characterized Class)",
      Matched_Count > 1  ~ "Level 3 (Ambiguous Isomeric Set)"
    )
  )

cat("Final Structured Metabolite Annotation Report:\n")
print(as.data.frame(annotation_report))
\end{rchunk}

\begin{routput}
Final Structured Metabolite Annotation Report:
  Feature\_ID Observed\_Mass Retention\_Time Matched\_Count
1     FT\_001      180.0639           2.45             3
2     FT\_002      153.0784           1.82             1
3     FT\_003      312.1245           8.40             0
                         Candidates Formulas Min\_PPM\_Error
1 Glucose | Fructose | myo-Inositol  C6H12O6          2.83
2                          Dopamine C8H11NO2          3.79
3                         Unmatched     None            NA
                            MSI\_Confidence
1         Level 3 (Ambiguous Isomeric Set)
2 Level 3 (Putatively Characterized Class)
3           Level 4 (Unidentified Feature)
\end{routput}

% =============================================================================
\section{The Five Deadly Sins of Metabolite Annotation}
% =============================================================================

\begin{enumerate}
    \item \textbf{Equating Accurate Mass with Confirmed Identity}: Accurate mass matches specify molecular formulas, not structures. Claiming Level 1 identification without matching retention time against an authentic standard analyzed locally is an egregious reporting violation.
    \item \textbf{Assuming Protonated Ions for All Peaks}: Assuming every feature in positive mode is $[M+\text{H}]^+$ introduces a $21.98\;\text{Da}$ mass error if the feature is actually $[M+\text{Na}]^+$, searching reference databases for completely wrong chemical families.
    \item \textbf{Using Excessive Mass Tolerances ($>15$ ppm)}: Applying loose search tolerances produces hundreds of false positive candidates, defeating the purpose of high-resolution mass spectrometry. Tolerances must reflect instrument calibration ($\le 5$ ppm).
    \item \textbf{Treating Adduct Multiplets as Independent Biomarkers}: Failing to deconvolute co-eluting adducts ($[M+\text{Na}]^+$, $[M+\text{K}]^+$, water losses) duplicates single metabolites across multiple feature rows, artificially inflating multiple-testing penalties in differential testing.
    \item \textbf{Ignoring In-Source Fragmentation}: Neutral losses ($\Delta m = -18.0106\;\text{Da}$ for water, $-43.9898\;\text{Da}$ for $\text{CO}_2$) generated in the ESI source are frequently misidentified as distinct endogenous metabolites.
\end{enumerate}

% =============================================================================
\section{Self-Assessment Diagnostic Checklist}
% =============================================================================

\begin{table}[ht]
\centering
\caption{Self-Assessment Diagnostic Checklist: Metabolite Annotation}
\label{tab:diagnostic_checklist}
\small
\begin{tabularx}{\textwidth}{l p{5.5cm} X c}
\toprule
\textbf{Category} & \textbf{Diagnostic Check} & \textbf{Pass Criterion} & \textbf{Status} \\
\midrule
\textbf{Deconvolution} & Co-eluting features grouped into pseudo-spectra. & Peak shape correlation $r \ge 0.85$ and $\Delta t_R \le 0.02$ min. & [ ] \\
\addlinespace
\textbf{Adduct Rules} & Neutral mass inverted from identified adduct ion. & Multi-adduct consistency check ($[M+\text{H}]^+$ and $[M+\text{Na}]^+$ agree). & [ ] \\
\addlinespace
\textbf{Mass Tolerance} & Symmetric parts-per-million window applied. & Search window width $\tau \le 5.0$ ppm with lock-mass calibration. & [ ] \\
\addlinespace
\textbf{Isotope Pattern} & Experimental envelope compared to candidate formula. & Isotope dot-product score $S_{\text{iso}} \ge 0.90$. & [ ] \\
\addlinespace
\textbf{Fine Structure} & High-resolution mass shifts evaluated. & Resolving power $R \ge 24,000$ to resolve $^{34}\text{S}$ from $^{13}\text{C}_2$. & [ ] \\
\addlinespace
\textbf{MSI Reporting} & Annotations mapped to formal MSI tiers. & Level 1 reserved for authentic standards; Level 3 for unresolved isomers. & [ ] \\
\bottomrule
\end{tabularx}
\end{table}

\FloatBarrier

% =============================================================================
\section{Chapter Exercises}
% =============================================================================

\noindent\textbf{Exercise 9.1 (Adduct Arithmetic and Mass Inversion).}\enspace
The neurotransmitter dopamine has molecular formula $\text{C}_8\text{H}_{11}\text{NO}_2$.
\begin{enumerate}
    \item Using exact atomic masses ($^{12}\text{C} = 12.000000$, $^{1}\text{H} = 1.007825$, $^{14}\text{N} = 14.003074$, $^{16}\text{O} = 15.994915$), compute the exact monoisotopic neutral mass $M_0$.
    \item Derive the theoretical $m/z$ values for the following adduct ions: $[M+\text{H}]^+$, $[M+\text{Na}]^+$, $[M+\text{H}-\text{H}_2\text{O}]^+$, and $[M-\text{H}]^-$.
\end{enumerate}

\vspace{0.8em}
\noindent\textbf{Exercise 9.2 (Parts-Per-Million Mass Accuracy).}\enspace
An LC-MS experiment observes a chromatographic peak at $m/z = 341.1082$. A reference database lists protonated sucrose ($[M+\text{H}]^+$, $\text{C}_{12}\text{H}_{23}\text{O}_{11}^+$) at theoretical $m/z = 341.1078$.
\begin{enumerate}
    \item Compute the absolute mass deviation in Daltons and relative deviation in parts-per-million (ppm).
    \item Determine whether this feature matches sucrose under a strict $\pm 5.0\;\text{ppm}$ tolerance threshold.
\end{enumerate}

\vspace{0.8em}
\noindent\textbf{Exercise 9.3 (Isotopic Fine Structure Resolution).}\enspace
Consider an unknown ion observed at nominal mass $m/z = 257$. Candidate formula $A$ is protonated palmitic acid ($\text{C}_{16}\text{H}_{33}\text{O}_2^+$, theoretical $M+2$ driven by $^{13}\text{C}_2$). Candidate formula $B$ is a sulfur-containing peptide fragment ($\text{C}_{12}\text{H}_{21}\text{N}_2\text{O}_2\text{S}^+$, theoretical $M+2$ driven by $^{34}\text{S}$).
\begin{enumerate}
    \item Calculate the physical mass separation between the $^{34}\text{S}$ and $^{13}\text{C}_2$ peaks at $M+2$.
    \item Calculate the minimum instrument resolving power $R_{\text{FWHM}}$ required to achieve baseline separation between these two isobaric peaks.
\end{enumerate}

\vspace{0.8em}
\noindent\textbf{Exercise 9.4 (MSI Confidence Tier Classification).}\enspace
A clinical metabolomics study reports identifying 350 serum metabolites using single-stage Orbitrap full-scan MS by querying the Human Metabolome Database (HMDB) within a $5.0\;\text{ppm}$ tolerance window without acquiring tandem MS/MS spectra or running authentic standards.
\begin{enumerate}
    \item Classify the true MSI confidence level of these annotations according to Sumner et al. (2007).
    \item Identify the specific failure mode in reporting hexose sugars ($\text{C}_6\text{H}_{12}\text{O}_6$) in this study.
\end{enumerate}

% =============================================================================
\section{Worked-Out Solutions}
% =============================================================================

\subsection{Solution to Exercise 9.1}

\noindent\textbf{Part 1: Neutral Mass Derivation}\enspace
For dopamine ($\text{C}_8\text{H}_{11}\text{NO}_2$):
\begin{align}
M_0 &= 8(12.0000000) + 11(1.0078250) + 1(14.0030740) + 2(15.9949146) \nonumber \\
    &= 96.0000000 + 11.0860750 + 14.0030740 + 31.9898292 = 153.0789782 \approx 153.0790\;\text{Da}
\end{align}

\noindent\textbf{Part 2: Adduct Theoretical $m/z$}\enspace
Using electron-corrected proton and adduct mass shifts:
\begin{align}
[M+\text{H}]^+ &= 153.078978 + 1.007276 = 154.086254 \approx 154.0863\;\text{Da} \\
[M+\text{Na}]^+ &= 153.078978 + 22.989218 = 176.068196 \approx 176.0682\;\text{Da} \\
[M+\text{H}-\text{H}_2\text{O}]^+ &= 154.086254 - 18.010565 = 136.075689 \approx 136.0757\;\text{Da} \\
[M-\text{H}]^- &= 153.078978 - 1.007276 = 152.071702 \approx 152.0717\;\text{Da}
\end{align}

\subsection{Solution to Exercise 9.2}

\noindent\textbf{Part 1: Mass Deviation}\enspace
The absolute error in Daltons is:
\begin{equation}
\Delta m = m_{\text{obs}} - m_{\text{theo}} = 341.1082 - 341.1078 = +0.0004\;\text{Da}
\end{equation}
The relative mass error in parts-per-million (ppm) is:
\begin{equation}
\delta_{\text{ppm}} = \frac{\Delta m}{m_{\text{theo}}} \times 10^6 = \frac{0.0004}{341.1078} \times 10^6 = +1.1726 \approx +1.17\;\text{ppm}
\end{equation}

\noindent\textbf{Part 2: Tolerance Evaluation}\enspace
Because $|\delta_{\text{ppm}}| = 1.17\;\text{ppm} \le 5.0\;\text{ppm}$, the observed peak matches theoretical sucrose within the allowable search window.

\subsection{Solution to Exercise 9.3}

\noindent\textbf{Part 1: Mass Separation}\enspace
The mass shift for double $^{13}\text{C}$ substitution is:
\begin{equation}
\Delta m(^{13}\text{C}_2) = 2 \times 1.0033548 = 2.006710\;\text{Da}
\end{equation}
The mass shift for $^{34}\text{S}$ substitution is:
\begin{equation}
\Delta m(^{34}\text{S}) = 33.967867 - 31.972071 = 1.995796\;\text{Da}
\end{equation}
The physical separation is:
\begin{equation}
\Delta m = |2.006710 - 1.995796| = 0.010914\;\text{Da} = 10.914\;\text{mDa}
\end{equation}

\noindent\textbf{Part 2: Required Resolving Power}\enspace
At nominal $m/z = 257$, baseline separation requires:
\begin{equation}
R_{\text{FWHM}} \ge \frac{m}{\Delta m} = \frac{257}{0.010914} \approx 23,548
\end{equation}
An instrument operating at resolving power $R \ge 24,000$ resolves the lipid $^{13}\text{C}_2$ peak from the sulfur $^{34}\text{S}$ peak cleanly.

\subsection{Solution to Exercise 9.4}

\noindent\textbf{Part 1: MSI Level Classification}\enspace
Matching full-scan accurate mass against a database without chromatographic retention time matching against authentic reference standards or MS/MS fragmentation provides **MSI Level 4 (putative molecular formula)** or at best **MSI Level 3 (putatively characterized compound class)**. The authors' claim of "identifying" 350 metabolites conflates Level 4 with Level 1.

\noindent\textbf{Part 2: Hexose Isomer Failure}\enspace
For features matching formula $\text{C}_6\text{H}_{12}\text{O}_6$ ($m/z = 180.0634$), exact mass cannot determine whether the molecule is D-glucose, D-galactose, D-fructose, or myo-inositol. Reporting "Glucose" without chromatographic separation misrepresents an unresolved four-way stereoisomer set.

% =============================================================================
% References
% =============================================================================

\begin{thebibliography}{99}
\small

\bibitem{sumner2007proposed}
Sumner, L.W., Amberg, A., Barrett, D., et al.: Proposed minimum reporting standards for chemical analysis Chemical Analysis Working Group (CAWG) Metabolomics Standards Initiative (MSI).
\textit{Metabolomics} \textbf{3}(3), 211--221 (2007).
\doi{10.1007/s11306-007-0082-2}

\bibitem{schymanski2014identifying}
Schymanski, E.L., Jeon, J., Gulde, R., et al.: Identifying small molecules via high resolution mass spectrometry: Communicating confidence.
\textit{Environmental Science \& Technology} \textbf{48}(4), 2097--2098 (2014).
\doi{10.1021/es5002105}

\bibitem{kuhl2012camera}
Kuhl, C., Tautenhahn, R., B{\"o}ttcher, C., Larson, T.R., Neumann, S.: CAMERA: An integrated strategy for compound spectra extraction and annotation of liquid chromatography/mass spectrometry data sets.
\textit{Analytical Chemistry} \textbf{84}(1), 283--289 (2012).
\doi{10.1021/ac202450g}

\bibitem{kind2007seven}
Kind, T., Fiehn, O.: Seven Golden Rules for heuristic filtering of molecular formulas obtained by accurate mass spectrometry.
\textit{BMC Bioinformatics} \textbf{8}, 105 (2007).
\doi{10.1186/1471-2105-8-105}

\bibitem{rainer2022modular}
Rainer, J., Vicini, A., Salzer, L., et al.: A modular and expandable ecosystem for metabolomics data analysis in R.
\textit{Metabolites} \textbf{12}(2), 173 (2022).
\doi{10.3390/metabo12020173}

\bibitem{gross2017mass}
Gross, J.H.: Mass Spectrometry: A Textbook, 3rd edn.
Springer, Cham (2017).
\doi{10.1007/978-3-319-54398-7}

\bibitem{smith2006xcms}
Smith, C.A., Want, E.J., O'Maille, G., Abagyan, R., Siuzdak, G.: XCMS: processing mass spectrometry data for metabolite profiling using nonlinear peak alignment, matching, and identification.
\textit{Analytical Chemistry} \textbf{78}(3), 779--787 (2006).
\doi{10.1021/ac051437y}

\bibitem{tautenhahn2008centwave}
Tautenhahn, R., B{\"o}ttcher, C., Neumann, S.: Highly sensitive feature detection for high resolution LC/MS.
\textit{BMC Bioinformatics} \textbf{9}, 504 (2008).
\doi{10.1186/1471-2105-9-504}

\bibitem{wishart2022hmdb}
Wishart, D.S., Guo, A., Oler, E., et al.: HMDB 5.0: the Human Metabolome Database for 2022.
\textit{Nucleic Acids Research} \textbf{50}(D1), D622--D631 (2022).
\doi{10.1093/nar/gkab1062}

\bibitem{guijas2018metlin}
Guijas, C., Montenegro-Burke, J.R., Domingo-Almenara, X., et al.: METLIN: a technology platform for identifying knowns and unknowns.
\textit{Analytical Chemistry} \textbf{90}(5), 3156--3164 (2018).
\doi{10.1021/acs.analchem.7b04424}

\bibitem{horai2010massbank}
Horai, H., Arita, M., Kanaya, S., et al.: MassBank: a public repository for sharing mass spectral data for life sciences.
\textit{Journal of Mass Spectrometry} \textbf{45}(7), 703--714 (2010).
\doi{10.1002/jms.1777}

\bibitem{wang2016sharing}
Wang, M., Carver, J.J., Phelan, V.V., et al.: Sharing and community curation of mass spectrometry data with Global Natural Products Social Molecular Networking.
\textit{Nature Biotechnology} \textbf{34}(8), 828--837 (2016).
\doi{10.1038/nbt.3597}

\bibitem{duehrkop2019sirius}
D{\"u}hrkop, K., Fleischauer, M., Ludwig, M., et al.: SIRIUS 4: a rapid tool for turning tandem mass spectra into metabolite formulas and structures.
\textit{Nature Methods} \textbf{16}(4), 299--302 (2019).
\doi{10.1038/s41592-019-0344-8}

\bibitem{tsugawa2015msdial}
Tsugawa, H., Cajka, T., Kind, T., et al.: MS-DIAL: data-independent MS/MS deconvolution for comprehensive metabolome analysis.
\textit{Nature Methods} \textbf{12}(6), 523--526 (2015).
\doi{10.1038/nmeth.3393}

\bibitem{pluskal2010mzmine}
Pluskal, T., Castillo, S., Villar-Briones, A., Ore{\v{s}}i{\v{c}}, M.: MZmine 2: modular framework for processing, visualizing, and analyzing mass spectrometry-based molecular profile data.
\textit{BMC Bioinformatics} \textbf{11}, 395 (2010).
\doi{10.1186/1471-2105-11-395}

\bibitem{huber2015orchestrating}
Huber, W., Carey, V.J., Gentleman, R., et al.: Orchestrating high-throughput genomic analysis with Bioconductor.
\textit{Nature Methods} \textbf{12}(2), 115--121 (2015).
\doi{10.1038/nmeth.3252}

\bibitem{gatto2021msnbase}
Gatto, L., Rainer, J.: MSnbase: an R/Bioconductor package for isobaric and label-free mass spectrometry data handling, processing and visualisation.
\textit{Bioinformatics} \textbf{28}(2), 288--289 (2012).
\doi{10.1093/bioinformatics/btr645}

\bibitem{broadhurst2018guidelines}
Broadhurst, D., Goodacre, R., Reinke, S.N., et al.: Guidelines and considerations for the use of system suitability and quality control samples in mass spectrometry assays applied to untargeted clinical metabolomic studies.
\textit{Metabolomics} \textbf{14}(6), 72 (2018).
\doi{10.1007/s11306-018-1367-3}

\end{thebibliography}

\end{document}
